\FloatBarrier\clearpage
\subsection{Signal extraction from the observed data}
\label{sec:extractions}\label{sec:v5-extractions}

The extraction displays show the leading separated positive fluctuations in
each campaign, the common-coupling fits at 66 and 92~MeV, and the largest
negative fitted contributions. Selection uses the dense observed profiled
local statistic; it does not use the stress-centered global ordering. The
66~MeV common-coupling fit is the largest positive combined fluctuation,
with nominal local significance about 2.76. The second leading positive
feature of the full union is at 21~MeV and contains only 2015; 92~MeV is a
second useful multi-campaign comparison, containing 2016 and 2021.

The plotted local coordinate is
\begin{equation}
 r=\operatorname{sign}(\widehat A)
 \sqrt{2[\mathrm{NLL}(0)-\mathrm{NLL}(\widehat A)]}.
\end{equation}
It agrees with the usual nominal local excess significance on the positive
branch. Negative values describe an auxiliary deficit fit, with positive
Poisson expectations maintained, and do not represent a physical negative
coupling. The probability mappings and their different sign conventions
are defined in Section~\ref{sec:v5-global}.

\begin{table}[htbp]\centering\small
\begin{tabular}{lrrr}\toprule
Data & Leading excess & Next separated excess & Deepest deficit\\\midrule
2015 full & 51 (+3.14) & 21 (+2.52) & 19 ($-3.21$)\\
2016 full & 90 (+3.42) & 117 (+3.28) & 102 ($-4.50$)\\
2021 10\% & 78 (+2.81) & 65 (+2.40) & 71 ($-4.02$)\\
Combined union & 66 (+2.76) & 21 (+2.52) & 72 ($-3.49$)\\\bottomrule
\end{tabular}
\caption{Selected mass in MeV, followed by signed nominal local significance.
The 19~MeV deficit is at the 2015 search endpoint. These values rank selected
local fits and are not global significances. The additional combined
92~MeV fit has nominal local significance 2.42.}
\end{table}

Displays retain native analysis bins and add sideband context. Residuals
subtract the same unprofiled GP mean throughout. Inside the likelihood window,
separate curves show the profiled-background displacement, total displacement
and signal. Lines join saved bin predictions, shown as counts per MeV.
Multi-campaign sums use exact shared whole-bin edges and have no independent fit.

\widefig{0.98}{../figures/v5_extraction_combined_66.pdf}{Common-coupling extraction at
66~MeV, with a shared fitted amplitude $\widehat\epsilon^2=4.55\times10^{-6}$
and nominal local significance 2.76. Each campaign retains its own yield
conversion and resolution. Bars show counting error only. Colored shading
in the residuals is the zero-centered GP constraint width; it is neither
fitted-background uncertainty nor total residual uncertainty.}
{fig:v5-extraction-combined66}
\widefig{0.98}{../figures/v5_extraction_combined_92.pdf}{The 92~MeV common-coupling fit
uses 2016 and 2021 and gives $\widehat\epsilon^2=2.97\times10^{-6}$ with nominal
local significance 2.42. The separate preferred rates differ, as shown at the
end of this results section. Bars show counting error only; residual shading
shows the zero-centered GP constraint width, not fitted-background or total
residual uncertainty.}
{fig:v5-extraction-combined92}
\widefig{0.98}{../figures/v5_extraction_combined_72.pdf}{The deepest combined observed
deficit at 72~MeV has signed local significance $-3.49$ and an auxiliary
amplitude $\widehat\epsilon^2=-4.88\times10^{-6}$. The negative signal template
describes missing counts relative to the fitted background. Bars show counting
error only; residual shading shows the zero-centered GP constraint width,
not fitted-background or total residual uncertainty.}
{fig:v5-extraction-combined72}
\fig{0.98}{../figures/v5_extraction_2015_peaks.pdf}{The two leading separated
2015 excesses at 51 and 21~MeV. The 21~MeV feature lies close to the search
endpoint and has strong background-model dependence in the traditional-fit
cross-check. Bars show counting error only; residual shading shows the
zero-centered GP constraint width, not fitted-background or total residual
uncertainty.}
{fig:v5-extraction-2015}
\fig{0.98}{../figures/v5_extraction_2016_peaks.pdf}{The leading separated 2016
excesses at 90 and 117~MeV. These fits retain the inherited 2016 numerical and
background-source qualifications. Bars show counting error only; residual
shading shows the zero-centered GP constraint width, not fitted-background
or total residual uncertainty.}
{fig:v5-extraction-2016}
\fig{0.98}{../figures/v5_extraction_2021_peaks.pdf}{The leading separated
excesses in the released 2021 10\% sample at 78 and 65~MeV. These masses
are fixed by the existing observed scan and are distinct from the combined
66~MeV coordinate. Bars show counting error only; residual shading shows
the zero-centered GP constraint width, not fitted-background or total
residual uncertainty.}
{fig:v5-extraction-2021}
\widefig{0.98}{../figures/v5_extraction_individual_deficits.pdf}{Individual deficit
extractions at 19, 102 and 71~MeV in 2015, 2016 and 2021. The first is a search
endpoint. Bars show counting error only; residual shading is the zero-centered
GP constraint width, not fitted-background or total residual uncertainty.
These are local background diagnostics. The joint deficit scan is retained
in Appendix~\ref{sec:v5-deficit}.}
{fig:v5-extraction-deficits}
\FloatBarrier

\subsection{Correlated fluctuations and signal-induced echoes}
\label{sec:echo}\label{sec:v5-echo}

A positive signal can affect a neighboring mass test because the exclusion
window moves. Counts protected from background training at the signal mass
can enter the sidebands of a nearby hypothesis. The GP prediction then
changes, and the neighboring signal fit can prefer a deficit. A negative
fitted contribution can therefore arise without removing any physical events
from the generating spectrum.

The deterministic 2021 replay in Figure~\ref{fig:v5-signal-echo} demonstrates this mechanism with saved positive
templates and the current dense solver. Relative to its own background-only
reconstruction, a 66~MeV injection changes the fitted local coordinate at
71 and 72~MeV by $-1.646$ and $-1.740$; a 78~MeV injection produces changes
$-1.825$ and $-1.968$. These changes are expressed in units of the nominal
local-significance coordinate. They are not separately calibrated
significances of an echo. The selected generating strengths are 17,142 events
at 66~MeV, 19,273 at 78~MeV and 36,373 for the 65+78~MeV pair; the comparison
is not an equal-strength experiment or a fitted two-particle hypothesis.

\fig{0.96}{../figures/v5_signal_echo_dense_replay.pdf}{Positive injected signals
can induce nearby fitted oscillations through the moving-mask background
regression. The injection-induced change is
$\Delta r(m)=r_m(B+S)-r_m(B)$ for the same smooth reference spectrum. The
65+78~MeV pair can induce a deep fitted deficit even though both generating
signal components are nonnegative. These are conditional response
illustrations; their agreement with one selected dip does not establish
the identity of the observed parent peak.}
{fig:v5-signal-echo}

Figure~\ref{fig:v501-correlations} places the three individual correlation
matrices beside the combined field. The off-diagonal positive and negative
lobes show that adjacent fitted extrema need not be independent. In the
combined panel, the response basis retains correlations where campaigns
overlap and gives zero covariance where no data are shared. These are
conditional covariance predictions; the independent complete-Poisson
validation is summarized in Section~\ref{sec:v5-global}.

\fig{0.99}{../figures/v501_profiled_correlations_2x2.pdf}{Profiled-response
correlation matrices under the archived generating backgrounds: 2015 full at
top left, 2016 full at top right, 2021 10\% at bottom left, and the combined
search at bottom right. All four panels use the same dimensionless color scale.
Positive and negative entries describe same-direction and opposite-direction
statistical fluctuations. Dotted lines in the combined panel mark changes in
the active datasets. The matrices quantify dependence; they do not assign a
probability to a selected observed echo or describe correlations of physical
signal production.}
{fig:v501-correlations}

A correlation slice can be expressed in the same local-significance units.
With $z_m=(r_m-a_m)/s_m$, conditioning a Gaussian field on $z_{m_0}=1$ gives
$E[r_m-a_m\mid z_{m_0}=1]=s_m R_{m,m_0}$. Figure~\ref{fig:v5-conditional-echo} uses fixed
anchors from the existing candidate catalogue. They illustrate the expected
accompanying fluctuation under the background covariance; they are distinct
from the positive-signal injection replay above.
\fig{0.98}{../figures/v5_conditional_echo_response.pdf}{Conditional changes in
the local-significance coordinate following a one-standard-deviation field
fluctuation at 51~MeV in 2015, 90~MeV in 2016 and 78~MeV in 2021. The curves
are $s_mR_{m,m_0}$, with each archived offset removed. Positive and negative
lobes show the scale and sign of correlated fluctuations. These are Gaussian
conditional means, not injected particles, probabilities of observed echoes,
or independent significance gains.}
{fig:v5-conditional-echo}

Figure~\ref{fig:v5-removal} gives a complementary observed-data check. Replacing
the selected 2021 candidate regions changes the fitted 71~MeV deficit from
$-4.02$ to $-0.72$, yet the remote observed scan retains 92\% of its original
standard deviation. The corresponding retained fractions in 2015 and 2016
are 67\% and 78\%. Some neighboring structures are coupled to the replaced
regions, while other variation remains. The replacements are conditional
interventions and cannot identify a particle or create a new null ensemble.

\fig{0.98}{../figures/v5_observed_candidate_removal.pdf}{Observed scans before and
after conditional replacement of the two selected candidate regions in each
year. The first and second regions are, respectively, 45--57 and
18.75--23.25~MeV in 2015 (centers 51 and 21~MeV); 81.75--98.25 and
105.75--128.25~MeV in 2016 (90 and 117~MeV); and 72.875--82.875 and
60.375--69.75~MeV in 2021 (78 and 65~MeV). Boundaries follow the stored
whole-bin selections around the $2.25\sigma_m$ masks. The envelope spans ten conditional replacements and is not a confidence
band. Changes demonstrate how the local fit responds to altered candidate
bins; the untouched remote variation is quantified in
Appendix~\ref{sec:candidate-removal}. All curves retain their existing tested
masses and the same nominal local-significance mapping.}
{fig:v5-removal}
\FloatBarrier
